\section{Summability, energy and power (2.31--2.36)}

\Prob{2.31\quad Show that the following sequences are absolutely summable}

\textbf{(a)} $x_a[n]=\dfrac{1}{4^{n}}\mu[n]$:
\[
\sum_{n=-\infty}^{\infty}|x_a[n]|=\sum_{n=0}^{\infty}\Bigl(\tfrac14\Bigr)^{n}
=\frac{1}{1-\tfrac14}=\frac43<\infty . \qquad\checkmark
\]

\textbf{(b)} $x_b[n]=\dfrac{1}{(n+2)(n+3)}\mu[n]$. Partial fractions give a telescoping sum:
\[
\frac{1}{(n+2)(n+3)}=\frac{1}{n+2}-\frac{1}{n+3},
\]
so for the partial sums
\[
\sum_{n=0}^{K}\frac{1}{(n+2)(n+3)}
=\Bigl(\frac12-\frac13\Bigr)+\Bigl(\frac13-\frac14\Bigr)+\cdots+\Bigl(\frac{1}{K+2}-\frac{1}{K+3}\Bigr)
=\frac12-\frac{1}{K+3}.
\]
All terms are positive, hence $\sum_n |x_b[n]|=\lim_{K\to\infty}\bigl(\tfrac12-\tfrac{1}{K+3}\bigr)$:
\[
\boxed{\;\sum_{n}|x_b[n]|=\tfrac12<\infty .\;}\qquad\checkmark
\]

\Prob{2.32\quad Absolutely summable $\Rightarrow$ finite energy, but not conversely}

\emph{Direction 1.} Let $S=\sum_n|x[n]|<\infty$. By Problem 2.23, $|x[n]|\le S$ for all $n$.
Therefore
\[
\mathcal{E}_x=\sum_{n}|x[n]|^{2}=\sum_{n}|x[n]|\cdot|x[n]|
\;\le\;S\sum_{n}|x[n]|=S^{2}<\infty .
\]
So every absolutely summable sequence has finite energy. (Equivalently
$\|x\|_2\le\|x\|_1$.)

\emph{Direction 2 (counterexample).} Take
\[
x[n]=\frac{1}{n}\,\mu[n-1].
\]
Its energy is finite,
\[
\mathcal{E}_x=\sum_{n=1}^{\infty}\frac{1}{n^{2}}=\frac{\pi^{2}}{6}\approx 1.6449<\infty,
\]
but its absolute sum is the divergent harmonic series,
\[
\sum_{n=1}^{\infty}\frac{1}{n}=\infty .
\]
Hence a finite-energy sequence need not be absolutely summable:
\[
\boxed{\;\ell_1\subsetneq\ell_2 .\;}
\]

\Prob{2.33\quad The square-summable $x_1[n]$ of Eq.\ (2.35) is not absolutely summable}

Equation (2.35) defines
\[
x_1[n]=\begin{cases}\dfrac1n, & n\ge 1,\\[4pt] 0, & n\le 0 .\end{cases}
\]
\emph{Square-summable:} $\displaystyle\sum_{n=1}^{\infty}\frac{1}{n^{2}}=\frac{\pi^{2}}{6}<\infty$.
(Bound it without knowing the exact value: $\frac{1}{n^2}<\frac{1}{n(n-1)}=\frac{1}{n-1}-\frac1n$
for $n\ge2$, whose sum telescopes to $1$; so $\mathcal{E}\le 2$.)

\emph{Not absolutely summable:} the harmonic series diverges. A quick proof by grouping,
\[
\sum_{n=1}^{2^{K}}\frac1n
=1+\frac12+\underbrace{\Bigl(\frac13+\frac14\Bigr)}_{>\frac12}
      +\underbrace{\Bigl(\frac15+\cdots+\frac18\Bigr)}_{>\frac12}+\cdots
\;>\;1+\frac{K}{2}\;\xrightarrow[K\to\infty]{}\;\infty .
\]
\[
\boxed{\;\sum_n|x_1[n]|=\infty\quad\text{while}\quad\sum_n|x_1[n]|^{2}=\tfrac{\pi^2}{6}<\infty .\;}
\]
This is precisely the counterexample used in Problem 2.32.

\Prob{2.34\quad $x_2[n]=\dfrac{\cos\omega_c n}{\pi n}\mu[n-1]$ is square-summable but not absolutely summable}

\emph{Square-summable.} Since $|\cos\omega_c n|\le 1$,
\[
\sum_{n=1}^{\infty}\Bigl|\frac{\cos\omega_c n}{\pi n}\Bigr|^{2}
\;\le\;\frac{1}{\pi^{2}}\sum_{n=1}^{\infty}\frac{1}{n^{2}}
=\frac{1}{\pi^{2}}\cdot\frac{\pi^{2}}{6}=\frac16<\infty . \qquad\checkmark
\]

\emph{Not absolutely summable} (for $0<\omega_c<\pi$). Use $|\cos\theta|\ge\cos^{2}\theta$
together with the identity $\cos^{2}\theta=\tfrac12(1+\cos 2\theta)$:
\[
\sum_{n=1}^{\infty}\frac{|\cos\omega_c n|}{\pi n}
\;\ge\;\frac{1}{\pi}\sum_{n=1}^{\infty}\frac{\cos^{2}\omega_c n}{n}
=\underbrace{\frac{1}{2\pi}\sum_{n=1}^{\infty}\frac{1}{n}}_{=\;\infty}
\;+\;\underbrace{\frac{1}{2\pi}\sum_{n=1}^{\infty}\frac{\cos 2\omega_c n}{n}}_{\text{finite}} .
\]
The first sum diverges (harmonic); the second converges by Dirichlet's test, since the partial
sums of $\cos 2\omega_c n$ are bounded whenever $2\omega_c\not\equiv 0 \pmod{2\pi}$ and $1/n$
decreases monotonically to $0$. (Its value is in fact $-\ln|2\sin\omega_c|$.) A divergent plus
a finite quantity diverges, so
\[
\boxed{\;\sum_{n}|x_2[n]|=\infty\quad\text{while}\quad \mathcal{E}_{x_2}\le \tfrac16<\infty .\;}
\]
This sequence is the (causal half of the) ideal lowpass impulse response --- the classical
example of an FIR-untruncatable, non-absolutely-summable but finite-energy signal.

\Prob{2.35\quad Compute the energy of the following sequences}

\textbf{(a)} $x_a[n]=A\alpha^{n}\mu[n]$, $|\alpha|<1$:
\[
\mathcal{E}_a=\sum_{n=0}^{\infty}\bigl|A\alpha^{n}\bigr|^{2}
=|A|^{2}\sum_{n=0}^{\infty}\bigl(|\alpha|^{2}\bigr)^{n}
\]
\[
\boxed{\;\mathcal{E}_a=\frac{|A|^{2}}{1-|\alpha|^{2}} .\;}
\]
(The geometric series converges exactly because $|\alpha|<1$; as $|\alpha|\to1$ the energy
blows up, as it must.)

\textbf{(b)} $x_b[n]=\dfrac{1}{n^{2}}\mu[n-1]$:
\[
\mathcal{E}_b=\sum_{n=1}^{\infty}\frac{1}{n^{4}}=\zeta(4)
\]
\[
\boxed{\;\mathcal{E}_b=\frac{\pi^{4}}{90}\approx 1.0823 .\;}
\]

\Prob{2.36\quad Average power and energy}

Recall $\mathcal{E}_x=\sum_n|x[n]|^{2}$ and
$\mathcal{P}_x=\lim_{K\to\infty}\frac{1}{2K+1}\sum_{n=-K}^{K}|x[n]|^{2}$. A signal with finite
nonzero energy always has $\mathcal{P}=0$; a signal with finite nonzero power always has
$\mathcal{E}=\infty$.

\textbf{(a)} $x_1[n]=(-1)^{n}$, so $|x_1[n]|^{2}=1$ for all $n$:
\[
\mathcal{E}_1=\sum_{n}1=\infty,\qquad
\mathcal{P}_1=\lim_{K\to\infty}\frac{2K+1}{2K+1}
\]
\[
\boxed{\;\mathcal{E}_1=\infty,\qquad \mathcal{P}_1=1 .\;}
\]

\textbf{(b)} $x_2[n]=\mu[n]$: $|x_2[n]|^{2}=1$ for $n\ge0$ and $0$ for $n<0$, so of the
$2K+1$ terms only $K+1$ are nonzero:
\[
\mathcal{E}_2=\infty,\qquad
\mathcal{P}_2=\lim_{K\to\infty}\frac{K+1}{2K+1}
\]
\[
\boxed{\;\mathcal{E}_2=\infty,\qquad \mathcal{P}_2=\tfrac12 .\;}
\]
(Only half of the time axis is "on", which halves the power relative to (a).)

\textbf{(c)} $x_3[n]=n\mu[n]$:
\[
\sum_{n=0}^{K}n^{2}=\frac{K(K+1)(2K+1)}{6}
\;\Longrightarrow\;
\mathcal{P}_3=\lim_{K\to\infty}\frac{1}{2K+1}\cdot\frac{K(K+1)(2K+1)}{6}
=\lim_{K\to\infty}\frac{K(K+1)}{6}=\infty .
\]
\[
\boxed{\;\mathcal{E}_3=\infty,\qquad \mathcal{P}_3=\infty .\;}
\]
$x_3[n]$ is \emph{neither} an energy signal nor a power signal.

\textbf{(d)} $x_4[n]=A_o e^{j\omega_o n}$: $|x_4[n]|^{2}=|A_o|^{2}$ for every $n$:
\[
\boxed{\;\mathcal{E}_4=\infty,\qquad \mathcal{P}_4=|A_o|^{2} .\;}
\]
(The power of a complex exponential does not depend on $\omega_o$ --- its magnitude is
constant.)

\textbf{(e)} $x_5[n]=A\cos\bigl(\tfrac{2\pi n}{M}+\phi\bigr)$ is periodic with period $M$
(for integer $M$), so the average power is the mean over one period:
\[
\mathcal{P}_5=\frac{1}{M}\sum_{n=0}^{M-1}A^{2}\cos^{2}\!\Bigl(\frac{2\pi n}{M}+\phi\Bigr)
=\frac{A^{2}}{2}+\frac{A^{2}}{2M}\sum_{n=0}^{M-1}\cos\!\Bigl(\frac{4\pi n}{M}+2\phi\Bigr).
\]
The remaining sum is a full period of a complex exponential of frequency $2\cdot\frac{2\pi}{M}$
and vanishes whenever $2/M\notin\mathbb{Z}$, i.e.\ for $M>2$:
\[
\boxed{\;\mathcal{E}_5=\infty,\qquad \mathcal{P}_5=\frac{A^{2}}{2}\quad (M>2).\;}
\]
Degenerate cases: for $M=1$, $x_5[n]=A\cos\phi$ and for $M=2$, $x_5[n]=A(-1)^{n}\cos\phi$;
in both, $\mathcal{P}_5=A^{2}\cos^{2}\phi$.
